\MFQTransition{Alpha 决策与研究材料答案}
\begin{enumerate}[leftmargin=*]
  \item MSE 服务数值预测，排序损失服务横截面次序，收益加权目标接近决策效用；三者仍需成本后组合检验。
  \item 零成交意味着订单没有改变旧仓位。只有成交的订单增量才进入新仓位。
  \item 第一期目标/实际为 $[1,0,-1]$；第二期目标 $[-1,1,0]$、订单 $[-2,1,1]$、成交后 $[0,1,-1]$。毛收益 $0.03+0.02=0.05$，换手 4、成本 0.004、净收益 0.046。
  \item 若 $L_R=-n^{-1}\sum s_ir_i$ 且没有范数或仓位约束，把正确方向分数乘任意正数会继续降低损失。
  \item 借券失败将对应卖单成交比例设为 0，并保留旧仓位；换手只累计实际成交增量。
  \item 首次发布时间晚于决策、抓取时使用未来版本、修订日晚于决策都应分别拒绝并给出稳定诊断。
  \item 好预测不等于可交易。应从分数重建目标、实际成交、成本和容量；净收益为负时不能以预测指标替代结论。
  \item 合格 Capstone 包含许可与版本、基线、checkpoint、嵌套时序、校准/漂移、文本审计、仓位成交、成本容量、故障注入、限制和一键命令。
\end{enumerate}

\section*{逐题推导与核对}
\begin{enumerate}[leftmargin=*]
  \item MSE 最小化数值预测误差；排序损失只关心横截面顺序；收益加权目标把预测映射到持仓和效用。要宣称可交易，还要继续通过时间、成交、成本、容量和漂移审计，不能把一个目标的改善当作所有目标改善。
  \item 目标仓位是意图，成交仓位是实际状态。未成交订单的增量为零，旧仓位必须保留；直接把未成交资产归零会制造虚假换手和收益。账本应同时保存 target、order、fill、position 四列。
  \item 第一期从 $[0,0,0]$ 到目标 $[1,0,-1]$，订单为 $[1,0,-1]$；第二期目标 $[-1,1,0]$，实际仓位从上一期为 $[1,0,-1]$ 出发，订单 $[-2,1,1]$，按成交比例 $[0.5,1,0]$ 得到 $[0,1,-1]$（第三只未成交，保留旧空仓）。按题设账本毛收益 $0.05$、换手 $4$、成本 $0.004$，净收益为 $0.046$；每一步都要写价格和持有期，不能只报总数。
  \item 若 $L_R=-n^{-1}\sum_i s_ir_i$ 没有分数范数或仓位约束，把任意正确方向的 $s$ 乘 $c>1$ 会使目标继续下降。加入 $\|s\|_2^2$、$\|s\|_1$、单名上限或风险预算后，优化才有界；这些约束的经济含义也要在报告中声明。
  \item 借券失败时卖单实际成交比例为零，保留旧空头/多头仓位；停牌同理。turnover 只累加实际 fill 的绝对增量，成本按实际成交而非目标订单计。把失败样本删掉会选择性保留容易成交的交易，必须作为显式状态进入统计。
  \item 三个泄漏负例分别是：决策时首次发布时间尚未发生；抓取版本后来被修订但回填到旧日；修订时间晚于决策却被当成可见正文。程序应按 \texttt{decision\_time >= available\_time} 断言并给不同失败码，不能把所有时间合并成一个日期。
  \item 审计预测优秀但成本后为负的报告：先重建分数到 target，再由 order/fill 得到实际仓位，最后逐日扣佣金、价差、冲击、借券和未成交机会成本。若净收益为负，应保留预测指标但把“可交易”结论改为不成立，不能用毛收益替代。
  \item Capstone 依次固定数据许可与截止、基线和 checkpoint、嵌套时间协议、文本/模型版本、预测到订单与成交、成本容量、漂移、失败注入和限制。最终包应能从原始输入重建每条 \texttt{target}、\texttt{order}、\texttt{fill}、\texttt{position} 与 \texttt{net\_return} 记录，并保存独立 oracle 和命令输出。
\end{enumerate}
